\section{Discussion and Threats to Validity}
\label{sec:threats}
\paragraph{Meaning of the guarantees.}
Safety and eventual progress are different obligations~\cite{safety}.
Our retry theorem fixes corrected intent, source membership, admission
predicates, and relevant channel inputs. Some stability lemmas are direct
consequences of that invariant. Their operational value is to expose the
conditions that a retry must preserve, not to claim that any repository
mutation is harmless. New commits, edits to release policy, or tag rewrites
can require a new plan. A process that fails forever need not converge.

The model abstracts external artifact contents and treats successful
publication followed by a completion record as one transition. Lost responses
and failed recording remain integration obligations. A Git lock coordinates
cooperating jobs but is not a transaction manager for arbitrary registries.
The implementation's own-reason failure exceptions, recording stages, and
finalization behavior require more detailed refinement than the abstract
blocked-set rule. We do not claim that testing selected traces closes that gap.

\paragraph{Internal and construct validity.}
The author developed the system, model, oracle, and tests, so common mistakes
can survive their agreement. The restricted oracle does not cover all current
semantics; full-profile categories are heuristic. A passing integrity check
shows consistency of supplied records, not independent replication. The
concurrent-push cases cover selected interleavings, not a systematic sweep.
The measured counts use different units: histories, package rows, episodes,
and assertions must not be pooled into one success percentage.

\paragraph{External validity.}
Synthetic repositories and local fault fixtures establish behavior under
controlled conditions. The historical subjects require reconstructed policies,
and the comparison uses pinned tool configurations. No user study establishes
that authors write correct intent or that teams find it easier to review.
Polyrepo applicability follows from the control-repository construction; its
operational scale is not measured. Cross-platform scripts, agent gates, and
multiple ecosystems demonstrate integration scope, not general effectiveness.
Performance observations from one contended host are not service-level promises.

\section{Related Work}
\label{sec:related}
\paragraph{Release and build systems.}
Adams and McIntosh identify release engineering as a distinct research
problem~\cite{adams}. Large monorepo practice~\cite{potvin} and build-system
models~\cite{mokhov} address dependency ordering and reuse. Incremental build
work also formalizes dynamic dependencies~\cite{pluto}; our model instead
uses a fixed release graph and records versioned publication obligations.
We claim neither incremental rebuilding nor a novel topological algorithm.
Lerna, Nx, Changesets, release-please, and semantic-release provide different
versioning and publication workflows~\cite{lerna,nx,changesets,releaseplease,
semanticrelease}. Section~\ref{sec:compare} compares operational traces of
specified configurations, not formal reductions of those complete tools.
Spinnaker already supports parallel pipeline stages and execution
histories~\cite{spinnaker}; expressing a DAG is not unique to this protocol.
Nix identifies component instances through input-derived store
paths~\cite{nix}, while TOSCA standardizes service templates and lifecycle
interfaces~\cite{tosca}. OpenGitOps calls for versioned desired state and
continuous reconciliation~\cite{gitops}. Our ledger instead tracks completed
versioned releases across independent publication targets. We neither replace
these systems' state models nor claim their guarantees for arbitrary scripts.

Version intent is not a proof of API compatibility. Studies of versioning
and ecosystem change~\cite{raemaekers,bogart,decan} motivate explicit policy
but do not show that a declared bump is correct. Package-installability work
such as EDOS~\cite{edos} models a different consistency question: a set of
installable packages rather than the remaining obligations of an interrupted
publication. Formal VCS work studies contents, patches, and merges
~\cite{swierstraloh,hpt,mimram,mrdt}; Git-mining work explains why history and
reference interpretation need care~\cite{birdgit}. Our ledger uses those
repository objects while leaving content and merge correctness outside its
proofs.

\paragraph{Recovery and formal assurance.}
Platform-independent deployment models already formalize lifecycle and
version compatibility: Liu and Smith's buildbox includes build, install,
ship, and update transitions~\cite{liusmith}. Aeolus studies component states,
dependencies, and configuration reachability~\cite{aeolus}. These are direct
antecedents to a common lifecycle abstraction. Our question is narrower:
preserving versioned publication obligations across partial executions.
The NEPHELE orchestration report also explicitly uses saga patterns for
distributed lifecycle operations~\cite{nephele}; applying sagas to operations
is therefore not itself a novelty claim.
Sagas~\cite{sagas} and the limits of transaction abstractions~\cite{gray}
provide the recovery setting. Independently committed and idempotently
processable operations~\cite{helland07,helland12} explain why uncertain
outcomes need explicit treatment. Our contribution specializes admission and
discharge to versioned consumers; it is not the invention of forward recovery.
Anvil~\cite{anvil} verifies reconciliation properties of Kubernetes controller
implementations, whereas our handwritten proofs establish a release model and
our tests sample implementation behavior. This is a significant difference
in assurance. Methodological cautions about fuzz evaluation~\cite{klees}
likewise motivate separating core agreement from unsupported generalization.

\section{Conclusion}
\label{sec:conclusion}
A release graph can cross build and publish boundaries. Local-artifact
consumers may proceed after a provider builds; registry-artifact consumers
must wait for publication. Correct ordering alone does not ensure recovery:
a failed consumer must retain its release obligation after successful
providers leave the plan. Consumer-based admission and separate train and
fresh-work windows make that obligation explicit. Under the stated retry
invariant, the model preserves versions, discharges recorded work, and
converges when attempts make progress. Released-binary experiments support
selected behaviors while exposing limits in oracle coverage and historical
policy replay. Mechanized proofs and a fuller implementation refinement,
including recording failures and randomized interleavings, remain open
work~\cite{lamport}.

\section*{Data Availability}
A curated local review bundle contains harness code, selected raw records,
per-file checksums, and an offline verifier. It reaggregates reported campaign
counts but omits binaries, complete Git histories, and some large dossiers;
it is not a complete executable replication package. The supplemental
mixed-chain experiment has separate code, event logs, and verdicts. Anonymous
distribution is being prepared. Identifying source and artifact links are
withheld in this working review draft; no public artifact availability is
claimed. An archival replication package is planned for publication.

\section*{AI Assistance Disclosure}
Most AI-assisted work used Claude Fable 5, with additional assistance from
Claude Opus 5, Claude Fable 5.1, Gemini, and OpenAI Codex. Across these tools,
assistance covered exposition, proof arguments, harness code, experiment
execution, and results summaries; it extended beyond language editing.
Measurements come from program executions, not language-model predictions.
The human author is responsible for the claims, proofs, citations, and
interpretation.
